\documentclass[10pt,a4paper]{article}
\usepackage[utf8]{inputenc}
\usepackage[T1]{fontenc}
\usepackage[margin=1.8cm,top=1.4cm,bottom=1.2cm]{geometry}
\usepackage{booktabs}
\usepackage{tabularx}
\usepackage{xcolor}
\usepackage{enumitem}
\usepackage{parskip}
\setlist{nosep,leftmargin=1.2em}
\setlength{\parskip}{4pt}
\emergencystretch=3em
\pagestyle{empty}
\definecolor{accent}{RGB}{0,84,140}
\newcommand{\sect}[1]{\vspace{6pt}{\color{accent}\large\bfseries #1}\par\vspace{1pt}}
\newcommand{\code}[1]{\texttt{\small #1}}
\newcommand{\keyfig}[2]{\begin{tabular}[t]{@{}c@{}}{\color{accent}\Large\bfseries #1}\\[-1pt]{\footnotesize #2}\end{tabular}}

\begin{document}

{\color{accent}\LARGE\bfseries AWI-ESM3 LR: estimated CMIP7 output volume}\\[2pt]
{\small TCo95 + FESOM 2.7 core3 + REcoM + LPJ-GUESS \textbar\ scaled from the measured HR test year \code{cli119} \textbar\ 14 September 2026}

\sect{Result}
With the HR variable set, LR physics and land produce roughly \textbf{41\,GB per year}. REcoM adds about
\textbf{10\,GB}, giving a baseline of \textbf{about 50\,GB per year}, roughly a ninth of HR. Fields that are
switched off for HR, and daily 3D biogeochemistry, can raise this to at most about 180\,GB per year.

\vspace{2pt}
\begin{center}
\keyfig{$\approx$\,41\,GB}{physics and land}\hfill
\keyfig{+\,10\,GB}{REcoM, monthly}\hfill
\keyfig{$\approx$\,50\,GB}{baseline per year}\hfill
\keyfig{5--6.5\,TB}{baseline per 100 years}\hfill
\keyfig{$\le$\,180\,GB}{per year, all options}
\end{center}

\sect{Physics and land: scaled from HR}
HR's cmorized volume was measured per grid on one full model year (529 files) and divided by the ratio of cells,
for 3D ocean fields also by the ratio of vertical levels. Sizes are compressed on-disk sizes.

\begin{tabularx}{\linewidth}{@{}Xrlr@{}}
\toprule
\textbf{grid} & \textbf{HR per year} & \textbf{HR cells $\rightarrow$ LR cells} & \textbf{LR per year} \\
\midrule
native atmosphere and land & 274\,GB & TCo319 421\,120 $\rightarrow$ TCo95 40\,320 ($\div$10.4) & 26.3\,GB \\
regridded atmosphere       &  68\,GB & 1296$\times$640 $\rightarrow$ 400$\times$192 ($\div$10.8) & 6.3\,GB \\
FESOM nodes, 2D            &  56\,GB & DARS 3\,146\,761 $\rightarrow$ core3 220\,509 ($\div$14.3) & 3.9\,GB \\
FESOM nodes, 3D            &  54\,GB & plus 57 $\rightarrow$ 48 levels ($\div$17.0) & 3.2\,GB \\
FESOM elements, 2D and 3D  &  18\,GB & 6\,229\,020 $\rightarrow$ 426\,403 ($\div$14.6 / $\div$17.4) & 1.1\,GB \\
fixed, decadal, climatology &  7\,GB & written once per run or decade & $<$1\,GB \\
\midrule
\textbf{total} & \textbf{470\,GB} & & \textbf{$\approx$\,41\,GB} \\
\bottomrule
\end{tabularx}

\vspace{2pt}
Coarse fields compress slightly worse per value and the fixed overhead per file does not shrink, so 40--55\,GB is
the realistic range. The file count stays at about 500 per year.

\sect{Ocean biogeochemistry: measured from a REcoM test run}
HR has no REcoM, so this part comes from \code{RECOM\_concdriven\_04} (model year 1900), which already runs on the
LR ocean. It writes 74 FESOM streams that HR does not have, 13.2\,GB per year, all monthly: 53 in 3D
(12.4\,GB) and 21 in 2D (0.8\,GB). Four of them (\code{bolus\_u/v/w}, \code{std\_dens\_DIVbolus}, 1.9\,GB) are
eddy-parameterisation diagnostics rather than biogeochemistry, which leaves \textbf{about 11\,GB of REcoM output}.
For monthly 3D ocean fields, cmorizing keeps 82--94\,\% of the raw size (measured on HR \code{temp} and
\code{salt}), so REcoM adds \textbf{about 10\,GB per year} if its output maps onto CMIP7 variables as written.

The CMIP7 data request lists 178 ocean biogeochemistry entries (106 variables): 151 monthly and 27 daily. Seven
of the daily ones are 3D (\code{detoc}, \code{dissoc}, \code{phycalc}, \code{phypico}, \code{zmeso}, \code{zmicro},
\code{zooc}). REcoM currently writes nothing daily; producing those seven would need daily 3D output and
would add about \textbf{50\,GB per year}, about 7.6\,GB per field.

\sect{Options that raise the total}
\begin{tabularx}{\linewidth}{@{}Xr@{}}
\toprule
\textbf{if produced for LR} & \textbf{added per year} \\
\midrule
6-hourly on 7 pressure levels (ta, ua, va, hus, zg), fixed in the LR XIOS configuration & +4\,GB \\
3-hourly on 6 pressure levels (ta, ua, va, hus, wap) & +8\,GB \\
6-hourly on model levels (ta, ua, va, hus, zg), 91 or 137 levels & +30 to 50\,GB \\
daily 3D biogeochemistry (seven fields above) & +50\,GB \\
eddy-induced transport (\code{msftmmpa}, GM is on in LR), ocean tendencies, monthly land and snow fields & $<$\,3\,GB \\
\bottomrule
\end{tabularx}

\vspace{2pt}
Pressure- and model-level estimates use HR's measured compressed bytes per stored value: about 1\,byte for
model-level and 1.7 to 2.7\,bytes for pressure-level fields.

\sect{Not covered}
\begin{itemize}
\item \textbf{Compute time.} HR needs 33 node-hours per model year, but that is dominated by I/O on the large
      grids and does not scale reliably to LR.
\item \textbf{Run lengths.} All numbers are per model year; multiply by the length of each LR experiment.
\item \textbf{Mapping of REcoM streams.} Not all 74 streams become CMIP7 variables one to one; some requested
      variables are derived (integrals, sums over plankton groups). The estimate assumes the volume stays similar.
\end{itemize}

{\footnotesize Sources: HR \code{/scratch/a/a270092/pycmor\_hr/cli119}; LR mesh \code{/work/ab0246/a270092/input/fesom2/core3};\\
LR atmosphere and REcoM \code{/work/bb1469/a270092/runtime/awiesm3-develop-cc/} (\code{cc\_tco95\_c}, \code{RECOM\_concdriven\_04})}

\end{document}
